\section{分块归约与内存访问}
\label{l15:sec:memory}

\subsection{数据分段与同步边界}

归约树相邻两层之间存在数据依赖：后一层要读取前一层写出的部分和。线程块内的屏障 \mbox{\code{__syncthreads()}} 使块内参与线程在该点会合，并使屏障前的共享内存或全局内存访问对块内后续访问正确可见。普通核函数中的这个屏障只作用于本块，不能等待其他线程块完成它们的树层。

因此，大输入通常按块分段。统一记总输入长度为 $N$，每块线程数为 $B$，块号为 $b$，块内线程号为 $t$。未粗化时每块处理 $2B$ 个元素，段首为 $s_b=2Bb$，块数为 $G=\lceil N/(2B)\rceil$。各块在自己的数据段内部完成归约，再把块和交给原子累加或后续核函数。这里将一次全局归约拆成块级局部归约，各段最终仍合成一个结果。通常所说的分段归约（segmented reduction）则为每个给定分段分别输出一个结果，输出含义与本讲任务不同。

局部树的所有线程都必须遵守同一层间同步结构。即使某个线程已经不再执行加法，只要它仍在这一固定块内算法中，就继续经过对应的块屏障。因此，\mbox{\code{__syncthreads()}} 放在活动线程的条件分支之后。把屏障移入 \mbox{\code{if (t < stride)}} 内部，会使仍在执行的同块线程走向不一致的同步位置。

下面先用完整数据段比较两种线程映射。两段原地代码均要求 $B$ 为二的幂、启动时恰好使用 $B$ 个线程，且输入缓冲区具有完整的 $2BG$ 个元素；若 $N<2BG$，应预先在缓冲区尾部补 $0$。它们会覆盖输入，并使用单独的块结果数组。之后的共享内存实现将显式检查 $N$，从而直接处理未补齐的尾块。

\subsection{递增与递减步长的线程映射}

第一种映射让线程 $t$ 从相邻的两个位置 $s_b+2t$ 和 $s_b+2t+1$ 开始。第一次合并后，部分和存放在偶数下标；下一次再隔开两个位置合并。步长 $s$ 按 $1,2,4,\ldots,B$ 递增，参与条件为 $t\bmod s=0$。固定的写入位置 $i=s_b+2t$ 与逐渐增大的第二操作数位置 $i+s$ 共同形成相邻子树的合并。

\noindent\begin{minipage}{\linewidth}
\begin{lstlisting}[caption={递增步长：在全局输入数组中原地合并},label={l15:lst:code1}]
template<int B>
__global__ void reduce_interleaved(float* data,
                                   float* partial) {
    unsigned int t = threadIdx.x;
    size_t base = size_t(2) * B * blockIdx.x;
    size_t i = base + 2 * t;
    for (unsigned int s = 1; s <= B; s *= 2) {
        if (t % s == 0) {
            data[i] += data[i + s];
        }
        __syncthreads();
    }
    if (t == 0) partial[blockIdx.x] = data[base];
}
\end{lstlisting}
\end{minipage}

当步长为 $s$ 的一层结束时，每个保留下来的位置保存连续 $2s$ 个原始元素之和。下一层把两个相邻的长度 $2s$ 的区间合并成长度 $4s$ 的区间，因此可以逐层推出最终段首保存全部 $2B$ 个元素之和。图\ref{l15:fig:increasing} 给出 $B=8$ 时的下标变化。

\begin{figure}[htbp]
\centering
\begin{tikzpicture}[x=0.80cm,y=0.88cm,>=Stealth,every node/.style={font=\footnotesize}]
\foreach \i in {0,...,15}{
  \filldraw[fill=SoftBlue,draw=UIUCBlue,line width=.35pt] (\i,0) rectangle +(1,.36);
  \node at (\i+.5,.62) {$\i$};}
\foreach \r/\s in {1/1,2/2,3/4,4/8}{
  \node[anchor=east,font=\small] at (-.2,-\r+.18) {步长 $\s$};
  \foreach \i in {0,...,15}{
    \pgfmathtruncatemacro{\active}{mod(\i,2*\s)==0 ? 1 : 0}
    \ifnum\active=1
      \filldraw[fill=SoftBlue,draw=UIUCBlue,line width=.35pt] (\i,-\r) rectangle +(1,.36);
      \draw[->,line width=.45pt] (\i+.5,-\r+1)--(\i+.5,-\r+.36);
      \draw[->,line width=.45pt] (\i+\s+.5,-\r+1)--(\i+.5,-\r+.36);
    \else
      \draw[draw=UIUCBlue,line width=.35pt] (\i,-\r) rectangle +(1,.36);
    \fi}}
\end{tikzpicture}

\caption{递增步长归约。蓝色位置保存仍会被后续层使用的部分和，存放位置逐层变得稀疏。}
\label{l15:fig:increasing}
\end{figure}

这套映射有两个相互独立的性能问题。首先，合并访问（coalesced access）关注\emph{同一条内存指令下，相邻线程访问的地址}。在第一层，线程 $0,1,2,\ldots$ 的第一操作数分别位于 $0,2,4,\ldots$，第二操作数分别位于 $1,3,5,\ldots$。每个线程读取的两个位置虽然相邻，每条读指令在整个线程束上的地址却带有空隙，通常需要更多内存事务才能服务相同数量的有效数据。

其次，一个线程束（warp）由连续的 $32$ 个线程组成。步长增大后，活动线程变为 $0,2,4,\ldots$，再变为 $0,4,8,\ldots$，使许多线程束内同时存在执行加法和跳过加法的线程。这种控制分歧（control divergence）降低了每条指令的有效通道利用率。直到较晚阶段，部分线程束才会整体停止算术工作。

第二种映射让线程 $t$ 首先合并 $s_b+t$ 与 $s_b+t+B$，即将整个数据段的左右两半对应相加。随后步长按 $B,B/2,\ldots,1$ 递减，活动线程始终是连续前缀 $0\leq t<s$。

\noindent\begin{minipage}{\linewidth}
\begin{lstlisting}[caption={递减步长：使活动线程和保存位置保持连续},label={l15:lst:code2}]
template<int B>
__global__ void reduce_contiguous(float* data,
                                  float* partial) {
    unsigned int t = threadIdx.x;
    size_t base = size_t(2) * B * blockIdx.x;
    size_t i = base + t;
    for (unsigned int s = B; s > 0; s /= 2) {
        if (t < s) {
            data[i] += data[i + s];
        }
        __syncthreads();
    }
    if (t == 0) partial[blockIdx.x] = data[base];
}
\end{lstlisting}
\end{minipage}

每一层把长度 $2s$ 的有效前缀缩短为长度 $s$ 的前缀。两半是互不重叠的部分和集合，合并后仍包含相同的一组原始输入，因而最终总和保持不变。完整的 $2B$ 元素归约一共有 $1+\log_2B$ 层；只有在首次 $2B\to B$ 的合并已经完成之后，剩余层数才是 $\log_2B$。

\begin{figure}[htbp]
\centering
\begin{tikzpicture}[x=0.80cm,y=0.88cm,>=Stealth,every node/.style={font=\footnotesize}]
\foreach \i in {0,...,15}{
  \filldraw[fill=SoftBlue,draw=UIUCBlue,line width=.35pt] (\i,0) rectangle +(1,.36);
  \node at (\i+.5,.62) {$\i$};}
\foreach \r/\s in {1/8,2/4,3/2,4/1}{
  \node[anchor=east,font=\small] at (-.2,-\r+.18) {步长 $\s$};
  \foreach \i in {0,...,15}{
    \ifnum\i<\s
      \filldraw[fill=SoftBlue,draw=UIUCBlue,line width=.35pt] (\i,-\r) rectangle +(1,.36);
      \draw[->,line width=.45pt] (\i+.5,-\r+1)--(\i+.5,-\r+.36);
      \draw[->,line width=.45pt] (\i+\s+.5,-\r+1)--(\i+.5,-\r+.36);
    \else
      \draw[draw=UIUCBlue,line width=.35pt] (\i,-\r) rectangle +(1,.36);
    \fi}}
\end{tikzpicture}

\caption{递减步长归约。每层把右半部分并入左半部分，有效部分和始终紧凑地位于前缀中。}
\label{l15:fig:decreasing}
\end{figure}

在相同的读取指令下，活动线程的地址现在依次是 $s_b+t$ 或 $s_b+t+s$，两组访问各自连续。只要起始地址满足相应对齐条件，这种映射有利于合并访存。在活动线程数至少为 $32$ 且为 $32$ 的倍数时，活动范围包含若干完整线程束，因此每个线程束内部的分支结果一致；到最后不足 $32$ 个线程工作时，分歧才集中出现在第一个线程束中。

例如 $B=128$ 时，两种算法各层的活动线程数都依次为 $128,64,32,16,8,4,2,1$，但含有活动线程的线程束数不同。表\ref{l15:tab:warp-activity} 只统计执行加法的线程束，未参与加法的线程仍会按代码经过块屏障。

\begin{table}[htbp]
\centering
\caption{每块 $128$ 个线程时，两种映射的算术活动分布}
\label{l15:tab:warp-activity}
\begin{tabular}{@{}lrrrrrrrr@{}}
\toprule
树层序号 & 0 & 1 & 2 & 3 & 4 & 5 & 6 & 7\\
\midrule
活动线程数 &128&64&32&16&8&4&2&1\\
递增步长的活动线程束数 &4&4&4&4&4&4&2&1\\
递减步长的活动线程束数 &4&2&1&1&1&1&1&1\\
\bottomrule
\end{tabular}
\end{table}

还可以把表中的活动线程束数换算成算术通道利用率。只考虑这些加法指令，将每个含有活动线程的线程束按 $32$ 个通道计费，则递增步长占用的通道槽位为 $32(4+4+4+4+4+4+2+1)=864$，递减步长为 $32(4+2+1+1+1+1+1+1)=384$。两者均完成 $256-1=255$ 次必要加法，因此利用率分别为 $255/864\approx29.5\%$ 和 $255/384\approx66.4\%$。这项指标量化了线程映射造成的有效通道比例，不计访存、屏障和调度，也不是 SM 占用率或实测加速比。

两种树具有相同的必要加法次数和依赖深度。性能差别来自这些操作如何分配给线程、每次访存服务多少有效数据，以及线程束内多少通道实际做了有用工作。因此，单看 $O(N)$ 工作量与 $O(\log N)$ 深度，还无法区分这两段代码的性能。

\subsection{在共享内存中保存中间结果}

归约过程会不断读写存放部分和的位置。虽然数值每层都在变化，这些中间存储位置会被反复使用，因此适合放入由同一线程块访问的共享内存（shared memory）。让全局内存只提供原始输入并接收块结果，就可以减少归约树内部对全局地址的读写，也能保留原输入以供后续计算使用。

首次加载可以直接完成 $2B\to B$ 的合并：线程 $t$ 读取两个原始值，在寄存器中相加，再把结果写入共享数组的第 $t$ 项。这样共享数组只需要 $B$ 项，而无需保存 $2B$ 个原始输入。所有槽位写入并同步后，再对这 $B$ 项运行递减步长归约。

\noindent\begin{minipage}{\linewidth}
\begin{lstlisting}[caption={带边界检查的共享内存归约：每块输出一个部分和},label={l15:lst:code3}]
template<int B>
__global__ void reduce_shared(const float* input,
                              size_t n, float* partial) {
    static_assert(B > 0 && (B & (B - 1)) == 0,
                  "B must be a power of two");
    __shared__ float values[B];
    unsigned int t = threadIdx.x;
    size_t i = size_t(2) * B * blockIdx.x + t;

    float a = (i < n) ? input[i] : 0.0f;
    float b = (i + B < n) ? input[i + B] : 0.0f;
    values[t] = a + b;
    __syncthreads();

    for (unsigned int s = B / 2; s > 0; s /= 2) {
        if (t < s) values[t] += values[t + s];
        __syncthreads();
    }
    if (t == 0) partial[blockIdx.x] = values[0];
}
\end{lstlisting}
\end{minipage}

该核函数使用一维网格和一维线程块，并要求实际启动线程数等于模板参数 $B$，同时符合设备的块大小限制。输出数组至少容纳 $G$ 项，且与输入区间不重叠；例如另一块尚未完成读取时，不能把紧凑块结果写回输入数组前部。共享数组中每个槽位都由对应线程初始化；两个输入地址独立进行边界判断，保证尾块中的有效输入不会被遗漏。

以 $B=128$、某一尾块仅剩 $150$ 个有效输入为例，线程 $0$ 至 $21$ 各读取两个有效值，线程 $22$ 至 $127$ 各读取一个有效值并配上一个 $0$。整个块仍初始化 $128$ 个共享槽位，并沿同一棵树归约。线程是否读到真实输入，与线程是否参与后续集体计算是两件事；补零使缺少输入的槽位仍具有定义良好的值，也保持了固定的同步结构。

初始化后的第一道屏障保证任何线程读取别人的共享槽位时，该槽位已经写好。循环中，每层的写入目标是前半个有效区间，第二操作数来自后半个有效区间；不同活跃线程的写入位置互不相同，且不会覆盖同层其他线程所需的第二操作数。层末屏障又保证后一层看到已经完成的前一层结果。这个不变量同时给出正确性依据和同步位置的依据。

对一个完整块，按源码中的标量操作计数，原地全局内存树有 $2B-1$ 次必要加法，每次涉及两次读和一次写；共享内存版本只从全局输入读取 $2B$ 个值，再向块结果数组写一个值。缓存、寄存器保留和内存事务合并会影响实际访存次数，因此这种源码级比较用于识别访问结构，不能直接当作显存总线上的字节测量。

若输入为四字节单精度数，全部 $N$ 个输入各读取一次、$G$ 个块结果各写入一次，对应 $4N+4G$ 字节的有效全局数据。所做加法与 $N$ 同阶，所以数据搬移、同步以及指令利用率容易主导耗时。共享内存优化解决了中间结果的存放问题，树后半段的块级同步和活动通道减少仍然存在。

在上面的直接写法中，每块执行一次初始化屏障和 $\log_2B$ 次循环屏障，共 $1+\log_2B$ 次。最后一层之后的屏障对只由线程 $0$ 读取自身结果而言可以省略，这里保留统一的循环形式，便于后续比较。进一步减少同步，需要利用线程束内部更小的通信范围。
